\chapter{دستگاه معادلات خطی، بردار و ماتریس}
\label{ch:linear-equations}

نقطه‌ی شروع جبر خطی یک سؤال ساده است: \emph{چند معادله‌ی خطی با چند مجهول داریم؛ جوابشان چیست؟} برای پاسخ دادن به این سؤال، دو ابزار اصلی این درس یعنی \textbf{بردار} و \textbf{ماتریس} را معرفی می‌کنیم و می‌بینیم که حل دستگاه چیزی جز چند «عمل سطری» ساده نیست؛ عمل‌هایی که خودشان را هم می‌شود با ضرب ماتریس‌ها نوشت.

\section{دستگاه معادلات خطی}

\subsection{معادله‌ی خطی چیست؟}
معادله‌ای مثل
\[
a_1x_1+a_2x_2+\dots+a_nx_n=b
\]
که در آن $a_1,\dots,a_n$ و $b$ عددهای ثابت و $x_1,\dots,x_n$ مجهول هستند، یک \term{معادله‌ی خطی}{Linear Equation} است. «خطی» یعنی هر مجهول فقط با توان یک و فقط ضرب در یک عدد ثابت ظاهر می‌شود؛ نه ضرب دو مجهول در هم، نه توان، نه تابعی مثل $\sin$.

\begin{example}[خطی یا غیرخطی؟]
\begin{itemize}
  \item $4x_1-5x_2+2=x_1$ خطی است (می‌شود آن را به شکل $3x_1-5x_2=-2$ نوشت).
  \item $x_2=2\sqrt{6}\,x_1+\pi x_3$ خطی است (ضرایب می‌توانند هر عدد حقیقی باشند).
  \item $x_1x_2=1$، \ $x_1^2+x_2=3$ و $\sin x_1+x_2=0$ خطی نیستند.
\end{itemize}
\end{example}

\subsection{دستگاه معادلات خطی}
\begin{definition}[دستگاه معادلات خطی]
فرض کنید $\F$ یک میدان باشد (مثلاً $\R$ یا $\C$). یک \term{دستگاه $m$ معادله‌ی خطی با $n$ مجهول}{System of Linear Equations} به این شکل است:
\[
\begin{cases}
A_{11}x_1+A_{12}x_2+\dots+A_{1n}x_n=y_1\\
A_{21}x_1+A_{22}x_2+\dots+A_{2n}x_n=y_2\\
\qquad\vdots\\
A_{m1}x_1+A_{m2}x_2+\dots+A_{mn}x_n=y_m
\end{cases}
\]
که در آن $A_{ij}$ و $y_k$ عضوهای $\F$ هستند. می‌خواهیم $n$ اسکالر $x_1,\dots,x_n$ از $\F$ پیدا کنیم که همه‌ی این معادله‌ها را با هم برقرار کنند.
\begin{itemize}
  \item اگر $y_1=y_2=\dots=y_m=0$ باشد، دستگاه را \term{همگن}{Homogeneous} می‌نامیم.
  \item یک \term{جواب}{Solution} دستگاه، برداری مثل $\vs=\begin{bmatrix}s_1\\ \vdots\\ s_n\end{bmatrix}$ است که با قرار دادن $x_1=s_1,\dots,x_n=s_n$ همه‌ی معادله‌ها درست شوند.
  \item مجموعه‌ی همه‌ی جواب‌ها را \term{مجموعه‌ی جواب}{Solution Set} می‌نامیم.
\end{itemize}
\end{definition}

\begin{note}
اندیس‌ها را این‌طور بخوانید: در $A_{ij}$ عدد اول ($i$) شماره‌ی \textbf{معادله} (سطر) و عدد دوم ($j$) شماره‌ی \textbf{مجهول} (ستون) است. پس $A_{23}$ ضریب $x_3$ در معادله‌ی دوم است.
\end{note}

\subsection{تعبیر هندسی: تصویر سطری}
دستگاه ساده‌ی زیر را در نظر بگیرید:
\[
\begin{cases}
x-2y=1\\
3x+2y=11
\end{cases}
\]
هر معادله در صفحه‌ی $xy$ یک \textbf{خط} است. جواب دستگاه نقطه‌ای است که روی هر دو خط باشد؛ یعنی \emph{محل برخورد} دو خط. با کمی محاسبه (که در ادامه می‌بینیم) جواب $x=3,\ y=1$ است.

\begin{center}
\begin{tikzpicture}[scale=0.62]
  \draw[step=1, gray!20] (-1,-2) grid (6,6);
  \draw[->] (-1,0)--(6.3,0) node[below] {$x$};
  \draw[->] (0,-2)--(0,6.3) node[left] {$y$};
  \draw[very thick, primary, domain=-1:6] plot (\x, {(\x-1)/2}) node[right] {$x-2y=1$};
  \draw[very thick, excolor, domain=1.5:4.4] plot (\x, {(11-3*\x)/2}) node[below right] {$3x+2y=11$};
  \fill[warncolor] (3,1) circle (4pt) node[above right=2pt] {$(3,1)$};
\end{tikzpicture}
\end{center}

\subsection{شکل ماتریسی: $A\vx=\vb$}
همین دستگاه را می‌توان به شکل فشرده‌ی «ماتریس ضرب در بردار» نوشت:
\[
\underbrace{\begin{bmatrix}1&-2\\3&2\end{bmatrix}}_{A}
\underbrace{\begin{bmatrix}x\\y\end{bmatrix}}_{\vx}
=
\underbrace{\begin{bmatrix}1\\11\end{bmatrix}}_{\vb}.
\]
\begin{itemize}
  \item به $A$ \term{ماتریس ضرایب}{Coefficient Matrix} می‌گوییم.
  \item اگر ستون $\vb$ را کنار $A$ بچسبانیم، \term{ماتریس افزوده}{Augmented Matrix} به دست می‌آید:
  \[
  \aug{A}{\vb}=\left[\begin{array}{cc|c}1&-2&1\\3&2&11\end{array}\right].
  \]
  ماتریس افزوده تمام اطلاعات دستگاه را در خود دارد؛ فقط اسم مجهول‌ها را حذف کرده‌ایم.
\end{itemize}
معنای دقیق ضرب ماتریس در بردار را در بخش \ref{sec:matmul} می‌بینیم.

\subsection{تعبیر هندسی دوم: تصویر ستونی}
همین دستگاه را می‌شود به شکل دیگری هم خواند: به صورت \textbf{ترکیبی از ستون‌های} $A$:
\[
x\begin{bmatrix}1\\3\end{bmatrix}+y\begin{bmatrix}-2\\2\end{bmatrix}=\begin{bmatrix}1\\11\end{bmatrix}.
\]
یعنی سؤال این است: \emph{چه مقدار از بردار ستون اول و چه مقدار از بردار ستون دوم را با هم جمع کنیم تا به $\vb$ برسیم؟} پاسخ: ۳ تا از اولی و ۱ تا از دومی.

\begin{center}
\begin{tikzpicture}[scale=0.42,>=Latex]
  \draw[step=1, gray!20] (-3,-1) grid (5,12);
  \draw[->] (-3,0)--(5.4,0) node[below] {\small $x_1$};
  \draw[->] (0,-1)--(0,12.4) node[left] {\small $x_2$};
  \draw[->, very thick, primary] (0,0)--(1,3) node[right] {\small $\begin{bmatrix}1\\3\end{bmatrix}$};
  \draw[->, very thick, primary!60] (0,0)--(3,9) node[right] {\small $3\begin{bmatrix}1\\3\end{bmatrix}$};
  \draw[->, very thick, excolor] (3,9)--(1,11);
  \draw[->, very thick, excolor!80] (0,0)--(-2,2) node[left] {\small $\begin{bmatrix}-2\\2\end{bmatrix}$};
  \fill[warncolor] (1,11) circle (5pt) node[above right] {\small $\vb=\begin{bmatrix}1\\11\end{bmatrix}$};
\end{tikzpicture}
\end{center}

\begin{intuition}
دو نگاه به یک دستگاه داریم:
\begin{itemize}
  \item \textbf{تصویر سطری:} هر معادله یک خط (یا صفحه) است و جواب، محل برخورد آن‌هاست.
  \item \textbf{تصویر ستونی:} سمت چپ ترکیبی از ستون‌های $A$ است و دنبال ضرایبی می‌گردیم که این ترکیب را برابر $\vb$ کند.
\end{itemize}
تصویر ستونی در ادامه‌ی درس (فضای ستونی، رتبه و …) بسیار مهم‌تر می‌شود.
\end{intuition}

\subsection{ایده‌ی حذف}
روش اصلی حل دستگاه این است: مضربی از یک معادله را از معادله‌ی دیگر کم کنیم تا یک مجهول \emph{حذف} شود.
\[
\begin{cases}
x-2y=1\\
3x+2y=11
\end{cases}
\quad\xrightarrow{\text{\fa{معادله‌ی دوم منهای ۳ برابر اولی}}}\quad
\begin{cases}
x-2y=1\\
\phantom{3x+{}}8y=8
\end{cases}
\]
حالا از معادله‌ی دوم $y=1$ و با جایگذاری در اولی $x=1+2y=3$. به این روش \term{جایگذاری پسرو}{Back Substitution} می‌گوییم. در شکل ماتریسی:
\[
\underbrace{\begin{bmatrix}1&-2\\0&8\end{bmatrix}}_{U}\begin{bmatrix}x\\y\end{bmatrix}=\underbrace{\begin{bmatrix}1\\8\end{bmatrix}}_{\vc}.
\]
ماتریس $A$ به یک ماتریس \term{بالامثلثی}{Upper Triangular} مثل $U$ تبدیل شد؛ ماتریسی که زیر قطر اصلی‌اش همه صفر است. این ایده در فصل ۲ به الگوریتم کاملی به نام \textbf{حذف گاوسی} تبدیل می‌شود.

\subsection{وجود و یکتایی جواب}
یک دستگاه خطی ممکن است:
\begin{itemize}
  \item \textbf{هیچ جوابی نداشته باشد}؛ در این صورت آن را \term{ناسازگار}{Inconsistent} می‌نامیم.
  \item \textbf{دقیقاً یک جواب داشته باشد}؛
  \item \textbf{بی‌شمار جواب داشته باشد}.
\end{itemize}
در دو حالت آخر دستگاه \term{سازگار}{Consistent} است. در صفحه، این سه حالت یعنی دو خط موازی، دو خط متقاطع و دو خط منطبق:

\begin{center}
\begin{tikzpicture}[scale=0.55]
  \begin{scope}
    \draw[->] (-0.5,0)--(3.5,0); \draw[->] (0,-0.5)--(0,3.2);
    \draw[thick, primary] (-0.3,0.3)--(2.7,3); \draw[thick, excolor] (0.5,-0.4)--(3.3,2.1);
    \node[below] at (1.5,-0.7) {\fa{بدون جواب}};
  \end{scope}
  \begin{scope}[xshift=6cm]
    \draw[->] (-0.5,0)--(3.5,0); \draw[->] (0,-0.5)--(0,3.2);
    \draw[thick, primary] (-0.3,0.3)--(3,3); \draw[thick, excolor] (-0.3,2.8)--(3,0.2);
    \fill[warncolor] (1.38,1.55) circle (3pt);
    \node[below] at (1.5,-0.7) {\fa{یک جواب}};
  \end{scope}
  \begin{scope}[xshift=12cm]
    \draw[->] (-0.5,0)--(3.5,0); \draw[->] (0,-0.5)--(0,3.2);
    \draw[line width=3pt, primary!40] (-0.3,0.3)--(2.7,3); \draw[thick, excolor, dashed] (-0.3,0.3)--(2.7,3);
    \node[below] at (1.5,-0.7) {\fa{بی‌شمار جواب}};
  \end{scope}
\end{tikzpicture}
\end{center}

آیا حالت دیگری مثلاً «دقیقاً دو جواب» هم ممکن است؟ قضیه‌ی زیر می‌گوید نه.

\begin{theorem}[تعداد جواب‌ها]\label{thm:num-solutions}
یک دستگاه معادلات خطی روی $\R$ (یا $\C$) یا جواب ندارد، یا دقیقاً یک جواب دارد، یا بی‌شمار جواب دارد.
\end{theorem}
\begin{proof}
کافی است نشان دهیم اگر دستگاه $A\vx=\vb$ دو جواب متفاوت $\vu\neq\vv$ داشته باشد، آن‌گاه بی‌شمار جواب دارد. برای هر عدد $t$ بردار
\[
\vx_t=\vu+t(\vv-\vu)
\]
را در نظر بگیرید. هر معادله‌ی دستگاه به شکل $a_1x_1+\dots+a_nx_n=b$ است. چون $\vu$ و $\vv$ هر دو جواب هستند:
\[
\sum_j a_j\big(u_j+t(v_j-u_j)\big)=\sum_j a_ju_j+t\Big(\sum_j a_jv_j-\sum_j a_ju_j\Big)=b+t(b-b)=b.
\]
پس $\vx_t$ برای هر $t$ جواب است. چون $\vv-\vu\neq\zero$، به ازای $t$های مختلف بردارهای $\vx_t$ متفاوت‌اند؛ پس بی‌شمار جواب داریم.
\end{proof}

در کل درس دنبال پاسخ دو سؤال بنیادی هستیم:
\begin{enumerate}
  \item \textbf{وجود:} آیا دستگاه سازگار است؟ یعنی دست‌کم یک جواب دارد؟
  \item \textbf{یکتایی:} اگر جواب دارد، آیا جوابش یکتاست؟
\end{enumerate}

\section{بردارها و ضرب داخلی}

\subsection{بردار و عمل‌های آن}
یک \term{بردار}{Vector} در $\R^n$ فهرست مرتبی از $n$ عدد حقیقی است که معمولاً به صورت \textbf{ستونی} نوشته می‌شود:
\[
\vv=\begin{bmatrix}v_1\\v_2\\ \vdots\\ v_n\end{bmatrix}\in\R^n .
\]
به $v_i$ \textbf{مؤلفه‌ی} $i$-ام بردار می‌گوییم. برای صرفه‌جویی در جا گاهی می‌نویسیم $\vv=(v_1,\dots,v_n)$؛ منظور همان بردار ستونی است. در $\R^2$ و $\R^3$ هر بردار را می‌توان به شکل یک \emph{پیکان} از مبدأ تا نقطه‌ی $(v_1,v_2)$ یا $(v_1,v_2,v_3)$ تصور کرد.

عمل‌های اصلی روی بردارها مؤلفه‌به‌مؤلفه انجام می‌شوند:
\[
\vu+\vv=\begin{bmatrix}u_1+v_1\\ \vdots\\ u_n+v_n\end{bmatrix},\qquad
\vu-\vv=\begin{bmatrix}u_1-v_1\\ \vdots\\ u_n-v_n\end{bmatrix},\qquad
c\vv=\begin{bmatrix}cv_1\\ \vdots\\ cv_n\end{bmatrix}.
\]
جمع دو بردار از نظر هندسی همان \textbf{قاعده‌ی متوازی‌الاضلاع} است و ضرب در اسکالر $c$ بردار را $c$ برابر بلندتر می‌کند (و اگر $c<0$ باشد جهتش را برعکس می‌کند).

\begin{center}
\begin{tikzpicture}[scale=0.85,>=Latex]
  \draw[->] (-0.3,0)--(5,0); \draw[->] (0,-0.3)--(0,3.4);
  \fill[primary!8] (0,0)--(3,1)--(4,3)--(1,2)--cycle;
  \draw[->, very thick, primary] (0,0)--(3,1) node[below right] {$\vu$};
  \draw[->, very thick, excolor] (0,0)--(1,2) node[above left] {$\vv$};
  \draw[dashed, gray] (3,1)--(4,3); \draw[dashed, gray] (1,2)--(4,3);
  \draw[->, very thick, thmcolor] (0,0)--(4,3) node[above right] {$\vu+\vv$};
  \begin{scope}[xshift=7.2cm]
    \draw[->] (-2.3,0)--(3,0); \draw[->] (0,-1.5)--(0,3.4);
    \draw[->, line width=4pt, primary!30] (0,0)--(2.4,1.6);
    \draw[->, very thick, primary] (0,0)--(1.2,0.8) node[above left] {$\vv$};
    \node[primary!70] at (2.75,1.95) {$2\vv$};
    \draw[->, very thick, warncolor] (0,0)--(-1.8,-1.2) node[below] {$-1.5\vv$};
  \end{scope}
\end{tikzpicture}
\end{center}

\subsection{یک کاربرد: کدگذاری داده‌های دسته‌ای}
در یادگیری ماشین گاهی با داده‌های \term{دسته‌ای}{Categorical} (غیرعددی) مثل رنگ یا کلمه سروکار داریم. برای این‌که الگوریتم‌ها بتوانند با آن‌ها کار کنند، باید آن‌ها را به بردار تبدیل کنیم. یک روش رایج \term{کدگذاری یک‌داغ}{One-Hot Encoding} است: به هر دسته برداری نسبت می‌دهیم که یک مؤلفه‌اش ۱ و بقیه صفرند.
\[
\text{red}=\begin{bmatrix}1\\0\\0\end{bmatrix},\quad
\text{green}=\begin{bmatrix}0\\1\\0\end{bmatrix},\quad
\text{blue}=\begin{bmatrix}0\\0\\1\end{bmatrix}.
\]
اگر زبان ما فقط چهار کلمه داشته باشد:
\[
\text{apple}=\begin{bmatrix}1\\0\\0\\0\end{bmatrix},\quad
\text{cat}=\begin{bmatrix}0\\1\\0\\0\end{bmatrix},\quad
\text{house}=\begin{bmatrix}0\\0\\1\\0\end{bmatrix},\quad
\text{tiger}=\begin{bmatrix}0\\0\\0\\1\end{bmatrix}.
\]
این بردارها همان \textbf{بردارهای پایه‌ای استاندارد} $\ve_1,\ve_2,\ve_3,\ve_4$ هستند. اما این روش دو عیب دارد که با ابزار بعدی (ضرب داخلی) بهتر می‌فهمیم:
\begin{itemize}
  \item بردارها بسیار \textbf{تُنُک} \enterm{Sparse} هستند: تقریباً همه‌ی مؤلفه‌ها صفرند و اگر واژگان ۵۰ هزار کلمه داشته باشد، هر بردار ۵۰ هزار مؤلفه دارد.
  \item هیچ دو کلمه‌ای \textbf{شبیه هم} نیستند: «گربه» همان‌قدر به «ببر» نزدیک است که به «خانه»!
\end{itemize}

\subsection{ضرب داخلی}
حاصل‌ضرب دو عدد، یک عدد است. اما حاصل «ضرب داخلی» دو بردار \emph{بردار نیست}؛ یک \textbf{عدد} است.

\begin{definition}[ضرب داخلی (ضرب نقطه‌ای)]
برای $\va,\vb\in\R^n$، \term{ضرب داخلی}{Inner/Dot/Scalar Product} آن‌ها عدد زیر است:
\[
\va\cdot\vb=a_1b_1+a_2b_2+\dots+a_nb_n=\sum_{i=1}^n a_ib_i .
\]
این عدد را با نمادهای $\va\T\vb$، $\inner{\va}{\vb}$ و $\langle\va|\vb\rangle$ هم نشان می‌دهند.
\end{definition}

مثلاً
\[
\begin{bmatrix}2\\1\end{bmatrix}\cdot\begin{bmatrix}1\\5\end{bmatrix}=2\times1+1\times5=7 .
\]

\begin{theorem}[ویژگی‌های ضرب داخلی]\label{thm:dot-props}
برای هر $\va,\vb,\vc\in\R^n$ و هر عدد $\alpha$:
\begin{enumerate}
  \item \textbf{تقارن:} $\va\cdot\vb=\vb\cdot\va$؛ یعنی $\va\T\vb=\vb\T\va$.
  \item \textbf{توزیع‌پذیری:} $\va\cdot(\vb+\vc)=\va\cdot\vb+\va\cdot\vc$.
  \item \textbf{بیرون آمدن اسکالر:} $(\alpha\va)\cdot\vb=\alpha(\va\cdot\vb)=\va\cdot(\alpha\vb)$.
  \item \textbf{مثبت بودن:} $\va\cdot\va\ge0$ و $\va\cdot\va=0$ فقط وقتی $\va=\zero$.
\end{enumerate}
\end{theorem}
\begin{proof}
همه مستقیم از تعریف و ویژگی‌های اعداد حقیقی به دست می‌آیند:
\begin{enumerate}
  \item $\sum a_ib_i=\sum b_ia_i$ چون ضرب اعداد جابه‌جایی است.
  \item $\sum a_i(b_i+c_i)=\sum a_ib_i+\sum a_ic_i$.
  \item $\sum (\alpha a_i)b_i=\alpha\sum a_ib_i=\sum a_i(\alpha b_i)$.
  \item $\va\cdot\va=\sum a_i^2$ مجموع چند عدد نامنفی است، پس نامنفی است؛ و مجموع چند عدد نامنفی فقط وقتی صفر است که همه‌ی آن‌ها صفر باشند، یعنی همه‌ی $a_i=0$.
\end{enumerate}
\end{proof}

\begin{note}[ترانهاده‌ی یک عدد]
چون $\va\T\vb$ یک عدد (ماتریس $1\times1$) است، ترانهاده‌اش با خودش برابر است. برای همین تساوی‌های زیر همه درست‌اند:
\[
\va\cdot\vb=\va\T\vb=(\va\T\vb)\T=\vb\T(\va\T)\T=\vb\T\va=\vb\cdot\va .
\]
(از قاعده‌ی $(XY)\T=Y\T X\T$ استفاده کردیم که در قضیه‌ی \ref{thm:transpose} ثابت می‌شود.)
\end{note}

\subsection{طول بردار}
ضرب داخلی یک بردار در خودش، \textbf{مربع طول} آن است:
\[
\va\T\va=\sum_{i=1}^n a_i^2=\norm{\va}^2,\qquad \norm{\va}=\sqrt{a_1^2+\dots+a_n^2}.
\]
به $\norm{\va}$ \term{طول}{Length} یا \term{نُرم}{Norm} بردار می‌گوییم. مثلاً
\[
\vv=\begin{bmatrix}4\\3\end{bmatrix}\ \Rightarrow\ \vv\cdot\vv=16+9=25,\qquad \norm{\vv}=5 .
\]
در $\R^2$ این همان قضیه‌ی فیثاغورس است: پیکان $(4,3)$ وتر یک مثلث قائم‌الزاویه با اضلاع ۴ و ۳ است.

\subsection{تعبیر هندسی ضرب داخلی}
\begin{theorem}[ضرب داخلی و زاویه]\label{thm:dot-angle}
اگر $\va,\vb$ دو بردار ناصفر در $\R^2$ یا $\R^3$ باشند و $\theta$ زاویه‌ی بین آن‌ها باشد، آن‌گاه
\[
\va\cdot\vb=\norm{\va}\,\norm{\vb}\cos\theta .
\]
\end{theorem}
\begin{proof}
مثلثی با رئوس مبدأ، نوک $\va$ و نوک $\vb$ را در نظر بگیرید. اضلاع آن $\norm{\va}$، $\norm{\vb}$ و $\norm{\va-\vb}$ هستند و زاویه‌ی بین دو ضلع اول $\theta$ است. طبق \textbf{قانون کسینوس‌ها}:
\[
\norm{\va-\vb}^2=\norm{\va}^2+\norm{\vb}^2-2\norm{\va}\norm{\vb}\cos\theta .
\]
از طرف دیگر با ویژگی‌های ضرب داخلی (قضیه‌ی \ref{thm:dot-props}):
\[
\norm{\va-\vb}^2=(\va-\vb)\cdot(\va-\vb)=\va\cdot\va-2\,\va\cdot\vb+\vb\cdot\vb=\norm{\va}^2-2\,\va\cdot\vb+\norm{\vb}^2 .
\]
با مقایسه‌ی دو رابطه، $-2\,\va\cdot\vb=-2\norm{\va}\norm{\vb}\cos\theta$ و حکم ثابت می‌شود.
\end{proof}

\begin{intuition}[سایه]
اگر $\vb$ بردار یکه (طول ۱) باشد، $\va\cdot\vb=\norm{\va}\cos\theta$ دقیقاً طول (علامت‌دار) \textbf{سایه‌ی} $\va$ روی امتداد $\vb$ است؛ مثل سایه‌ی یک چوب وقتی نور عمود بر $\vb$ بتابد.
\begin{center}
\begin{tikzpicture}[scale=0.9,>=Latex]
  \draw[->, very thick, primary] (0,0)--(5,0) node[below] {$\vb$};
  \draw[->, very thick, excolor] (0,0)--(3,2.2) node[above] {$\va$};
  \draw[dashed] (3,2.2)--(3,0);
  \draw (3,0) rectangle (2.8,0.2);
  \draw[line width=3pt, thmcolor, opacity=0.6] (0,0)--(3,0);
  \node[below, thmcolor] at (1.5,-0.1) {\small $\norm{\va}\cos\theta$};
  \draw (0.8,0) arc (0:36:0.8); \node at (1.1,0.32) {\small $\theta$};
\end{tikzpicture}
\end{center}
پس ضرب داخلی نشان می‌دهد دو بردار چقدر «هم‌جهت» یا \textbf{شبیه} هم هستند:
\begin{itemize}
  \item اگر $\va\cdot\vb>0$ زاویه حاده است (هم‌جهت‌اند)؛
  \item اگر $\va\cdot\vb=0$ بر هم \textbf{عمود} \enterm{Orthogonal} هستند؛
  \item اگر $\va\cdot\vb<0$ زاویه منفرجه است (تقریباً خلاف جهت هم‌اند).
\end{itemize}
\end{intuition}

\begin{example}[عیب کدگذاری یک‌داغ]
نشان دهید هیچ دو کلمه‌ی متفاوتی در کدگذاری یک‌داغ «شبیه» هم نیستند.
\end{example}
\begin{solution}
\[
\text{apple}\cdot\text{cat}=
\begin{bmatrix}1\\0\\0\\0\end{bmatrix}\cdot\begin{bmatrix}0\\1\\0\\0\end{bmatrix}=0=
\begin{bmatrix}0\\0\\0\\1\end{bmatrix}\cdot\begin{bmatrix}0\\1\\0\\0\end{bmatrix}=\text{tiger}\cdot\text{cat}.
\]
در حالت کلی $\ve_i\cdot\ve_j=0$ برای $i\neq j$؛ یعنی همه‌ی کلمه‌ها بر هم عمودند و شباهت «گربه» به «ببر» با شباهتش به «سیب» یکی است. برای همین در عمل از بردارهای چگال‌تری (\en{word embeddings}) استفاده می‌کنند که در آن کلمه‌های هم‌معنا ضرب داخلی بزرگ‌تری دارند.
\end{solution}

\section{ماتریس‌ها و ضرب آن‌ها}\label{sec:matmul}

\subsection{نمادگذاری}
منظور از $A\in\R^{m\times n}$ ماتریسی با $m$ سطر و $n$ ستون است که درایه‌هایش اعداد حقیقی‌اند. یک ماتریس را می‌توانیم به سه شکل ببینیم: درایه‌به‌درایه، ستون‌به‌ستون یا سطربه‌سطر:
\[
A=\begin{bmatrix}
a_{11}&a_{12}&\cdots&a_{1n}\\
a_{21}&a_{22}&\cdots&a_{2n}\\
\vdots&\vdots&\ddots&\vdots\\
a_{m1}&a_{m2}&\cdots&a_{mn}
\end{bmatrix}
=
\begin{bmatrix}
| & | & & |\\
\va_1 & \va_2 & \cdots & \va_n\\
| & | & & |
\end{bmatrix}
=
\begin{bmatrix}
- & \va_1\T & -\\
- & \va_2\T & -\\
 & \vdots & \\
- & \va_m\T & -
\end{bmatrix}.
\]
(دقت کنید که در نمایش ستونی $\va_j\in\R^m$ ستون $j$-ام است و در نمایش سطری $\va_i\T$ سطر $i$-ام؛ از متن معلوم می‌شود منظور کدام است.)

\begin{definition}[چند ماتریس خاص]
\begin{itemize}
  \item \textbf{مربعی:} تعداد سطر و ستون برابر است ($m=n$).
  \item \textbf{صفر} ($0$): همه‌ی درایه‌ها صفرند.
  \item \textbf{همانی} ($I$ یا $I_n$): ماتریس مربعی که روی قطر اصلی ۱ و بقیه‌ی جاها ۰ است؛ یعنی $I_{ij}=\delta_{ij}$.
  \item \textbf{قطری:} مربعی که بیرون قطر اصلی صفر است.
  \item \textbf{بالامثلثی / پایین‌مثلثی:} مربعی که زیر / بالای قطر اصلی صفر است.
\end{itemize}
\end{definition}

\begin{definition}[جمع، ضرب اسکالر و ترانهاده]
برای $A,B\in\R^{m\times n}$ و عدد $c$:
\[
(A+B)_{ij}=A_{ij}+B_{ij},\qquad (cA)_{ij}=c\,A_{ij},\qquad (A\T)_{ij}=A_{ji}.
\]
ماتریس $A\T\in\R^{n\times m}$ را \term{ترانهاده}{Transpose} $A$ می‌نامیم؛ سطرهای $A$ ستون‌های $A\T$ می‌شوند.
\end{definition}

\subsection{ترکیب خطی}
\begin{definition}[ترکیب خطی]
یک \term{ترکیب خطی}{Linear Combination} از بردارهای $\va_1,\dots,\va_m$ (همه با اندازه‌ی $n$) هر برداری به شکل
\[
\beta_1\va_1+\beta_2\va_2+\dots+\beta_m\va_m
\]
است که $\beta_1,\dots,\beta_m$ اسکالر هستند. به $\beta_i$ها \textbf{ضرایب} ترکیب خطی می‌گوییم.
\end{definition}
مفهوم ترکیب خطی قلب جبر خطی است؛ تقریباً هر مفهوم بعدی (پوشش، استقلال، پایه، رتبه) با آن تعریف می‌شود.

\subsection{ضرب ماتریس در بردار}
\begin{definition}[ضرب ماتریس در بردار]
برای $A\in\R^{m\times n}$ و $\vx\in\R^n$، حاصل $\vy=A\vx$ برداری در $\R^m$ است با مؤلفه‌های
\[
y_i=\sum_{j=1}^n a_{ij}x_j,\qquad i=1,\dots,m .
\]
\end{definition}

این تعریف را می‌شود به دو شکل معادل دید:

\paragraph{نگاه سطری.} مؤلفه‌ی $i$-ام حاصل، ضرب داخلی سطر $i$-ام $A$ در $\vx$ است:
\[
A\vx=\begin{bmatrix}- & \va_1\T & -\\ - & \va_2\T & -\\ & \vdots & \\ - & \va_m\T & -\end{bmatrix}\vx
=\begin{bmatrix}\va_1\T\vx\\ \va_2\T\vx\\ \vdots\\ \va_m\T\vx\end{bmatrix}.
\]

\paragraph{نگاه ستونی.} حاصل، ترکیب خطی ستون‌های $A$ با ضرایب $x_1,\dots,x_n$ است:
\[
A\vx=\begin{bmatrix}| & | & & |\\ \va_1 & \va_2 & \cdots & \va_n\\ | & | & & |\end{bmatrix}\begin{bmatrix}x_1\\x_2\\ \vdots\\ x_n\end{bmatrix}
=x_1\va_1+x_2\va_2+\dots+x_n\va_n .
\]

\begin{proof}[هم‌ارزی دو نگاه]
مؤلفه‌ی $i$-ام بردار $x_1\va_1+\dots+x_n\va_n$ برابر است با $x_1a_{i1}+\dots+x_na_{in}=\sum_j a_{ij}x_j$، که همان $y_i$ و همان $\va_i\T\vx$ (ضرب داخلی سطر $i$ در $\vx$) است.
\end{proof}

\begin{note}
نگاه ستونی می‌گوید: \hl{$A\vx$ همیشه ترکیبی خطی از ستون‌های $A$ است.} به‌زودی خواهیم دید که «ستون‌های $A$ مستقل خطی‌اند اگر و تنها اگر از $A\vx=\zero$ نتیجه شود $\vx=\zero$».
\end{note}

\paragraph{ضرب بردار سطری در ماتریس.} می‌توانیم از سمت چپ هم یک بردار سطری $\vx\T$ (با $\vx\in\R^m$) را در $A$ ضرب کنیم. باز دو نگاه داریم:
\[
\vy\T=\vx\T A=\vx\T\begin{bmatrix}| & & |\\ \va_1 & \cdots & \va_n\\ | & & |\end{bmatrix}=\begin{bmatrix}\vx\T\va_1 & \vx\T\va_2 & \cdots & \vx\T\va_n\end{bmatrix},
\]
و با نوشتن $A$ بر حسب سطرها:
\[
\vx\T A=\begin{bmatrix}x_1 & x_2 & \cdots & x_m\end{bmatrix}\begin{bmatrix}- & \va_1\T & -\\ & \vdots & \\ - & \va_m\T & -\end{bmatrix}
=x_1\big[-\ \va_1\T\ -\big]+x_2\big[-\ \va_2\T\ -\big]+\dots+x_m\big[-\ \va_m\T\ -\big].
\]
یعنی \hl{$\vx\T A$ ترکیبی خطی از سطرهای $A$ است.}

\begin{theorem}[ویژگی‌های ضرب ماتریس در بردار]\label{thm:matvec-props}
برای ماتریس‌های $A,B\in\R^{m\times n}$، بردارهای $\vu,\vv\in\R^n$ و عدد $\alpha$:
\[
A(\vu+\vv)=A\vu+A\vv,\quad (A+B)\vu=A\vu+B\vu,\quad (\alpha A)\vu=\alpha(A\vu)=A(\alpha\vu),
\]
\[
0\,\vu=\zero,\qquad A\zero=\zero,\qquad I\vu=\vu .
\]
\end{theorem}
\begin{proof}
همه را مؤلفه‌به‌مؤلفه بررسی می‌کنیم. مثلاً مؤلفه‌ی $i$-ام $A(\vu+\vv)$ برابر است با
\[
\sum_j a_{ij}(u_j+v_j)=\sum_j a_{ij}u_j+\sum_j a_{ij}v_j ,
\]
که مؤلفه‌ی $i$-ام $A\vu+A\vv$ است. بقیه هم به همین سادگی‌اند؛ برای آخری: مؤلفه‌ی $i$-ام $I\vu$ برابر $\sum_j\delta_{ij}u_j=u_i$ است چون در این جمع فقط جمله‌ی $j=i$ ناصفر است.
\end{proof}

\begin{example}
برای ماتریس
$A=\begin{bmatrix}-1&2&1\\2&0&-2\end{bmatrix}$
موارد زیر را به صورت «ضرب ماتریس و بردار» بنویسید و حساب کنید: ستون $j$-ام، سطر $i$-ام، مجموع سطرها و مجموع ستون‌ها.
\end{example}
\begin{solution}
از بردارهای پایه‌ای استاندارد $\ve_j$ (بردار با ۱ در جای $j$ و صفر در بقیه) و بردار تمام‌یک $\mathbf{1}$ استفاده می‌کنیم:
\begin{itemize}
  \item \textbf{ستون $j$-ام:} $\va_j=A\ve_j$؛ چون طبق نگاه ستونی $A\ve_j=0\va_1+\dots+1\va_j+\dots+0\va_n$. مثلاً $A\ve_2=\begin{bmatrix}2\\0\end{bmatrix}$.
  \item \textbf{سطر $i$-ام:} $\va_i\T=\ve_i\T A$. مثلاً $\ve_1\T A=\begin{bmatrix}-1&2&1\end{bmatrix}$.
  \item \textbf{مجموع سطرها:} $\mathbf{1}\T A=\begin{bmatrix}1&1\end{bmatrix}A=\begin{bmatrix}1&2&-1\end{bmatrix}$.
  \item \textbf{مجموع ستون‌ها:} $A\mathbf{1}=A\begin{bmatrix}1\\1\\1\end{bmatrix}=\begin{bmatrix}2\\0\end{bmatrix}$.
\end{itemize}
\end{solution}

\subsection{ضرب ماتریس در ماتریس}
\begin{definition}[ضرب دو ماتریس]
فرض کنید $A$ یک ماتریس $m\times n$ و $B$ یک ماتریس $n\times p$ روی میدان $\F$ باشند. حاصل‌ضرب $AB$ ماتریس $m\times p$ مثل $C$ است که درایه‌ی $(i,j)$ آن
\[
C_{ij}=\sum_{r=1}^{n}A_{ir}B_{rj}
\]
است. به بیان دیگر $C_{ij}=\va_i\T\vb_j$: ضرب داخلی \textbf{سطر $i$-ام $A$} در \textbf{ستون $j$-ام $B$}.
\end{definition}

\begin{warning}[شرط ابعاد]
ضرب $AB$ فقط وقتی تعریف شده است که \textbf{تعداد ستون‌های $A$} با \textbf{تعداد سطرهای $B$} برابر باشد:
\[
\underbrace{A}_{m\times n}\ \underbrace{B}_{n\times p}=\underbrace{C}_{m\times p}.
\]
\end{warning}

\begin{example}
$\begin{bmatrix}1&2\\3&4\end{bmatrix}\begin{bmatrix}5&6\\7&8\end{bmatrix}$ را حساب کنید.
\end{example}
\begin{solution}
\[
\begin{bmatrix}1\cdot5+2\cdot7 & 1\cdot6+2\cdot8\\ 3\cdot5+4\cdot7 & 3\cdot6+4\cdot8\end{bmatrix}=\begin{bmatrix}19&22\\43&50\end{bmatrix}.
\]
\end{solution}

\subsubsection{چهار نگاه به ضرب ماتریس‌ها}
فرض کنید $A\in\R^{m\times n}$ با سطرهای $\va_i\T$ و ستون‌های $\va_j$، و $B\in\R^{n\times p}$ با سطرهای $\vb_i\T$ و ستون‌های $\vb_j$ باشد.

\paragraph{۱. مجموعه‌ای از ضرب‌های داخلی} (همان تعریف):
\[
C=AB=\begin{bmatrix}- & \va_1\T & -\\ - & \va_2\T & -\\ & \vdots & \\ - & \va_m\T & -\end{bmatrix}
\begin{bmatrix}| & | & & |\\ \vb_1 & \vb_2 & \cdots & \vb_p\\ | & | & & |\end{bmatrix}
=\begin{bmatrix}
\va_1\T\vb_1 & \va_1\T\vb_2 & \cdots & \va_1\T\vb_p\\
\va_2\T\vb_1 & \va_2\T\vb_2 & \cdots & \va_2\T\vb_p\\
\vdots & \vdots & \ddots & \vdots\\
\va_m\T\vb_1 & \va_m\T\vb_2 & \cdots & \va_m\T\vb_p
\end{bmatrix}.
\]

\paragraph{۲. مجموع ضرب‌های خارجی.} برای $\vu\in\R^m$ و $\vv\in\R^p$، ماتریس $m\times p$ مثل $\vu\vv\T$ را \term{ضرب خارجی}{Outer Product} می‌نامیم؛ درایه‌ی $(i,j)$ آن $u_iv_j$ است. حالا:
\[
C=AB=\begin{bmatrix}| & | & & |\\ \va_1 & \va_2 & \cdots & \va_n\\ | & | & & |\end{bmatrix}
\begin{bmatrix}- & \vb_1\T & -\\ - & \vb_2\T & -\\ & \vdots & \\ - & \vb_n\T & -\end{bmatrix}
=\sum_{r=1}^{n}\va_r\vb_r\T .
\]
\textbf{چرا؟} درایه‌ی $(i,j)$ ماتریس $\va_r\vb_r\T$ برابر $A_{ir}B_{rj}$ است؛ پس درایه‌ی $(i,j)$ مجموع برابر $\sum_r A_{ir}B_{rj}=C_{ij}$.

\paragraph{۳. مجموعه‌ای از ضرب‌های ماتریس در بردار.} ستون $j$-ام $C$ برابر $A\vb_j$ است:
\[
C=AB=A\begin{bmatrix}| & | & & |\\ \vb_1 & \vb_2 & \cdots & \vb_p\\ | & | & & |\end{bmatrix}=\begin{bmatrix}| & | & & |\\ A\vb_1 & A\vb_2 & \cdots & A\vb_p\\ | & | & & |\end{bmatrix}.
\]
\textbf{چرا؟} مؤلفه‌ی $i$-ام $A\vb_j$ برابر $\sum_r A_{ir}B_{rj}=C_{ij}$ است. پس طبق نگاه ستونی ضرب ماتریس در بردار، \hl{هر ستون $AB$ ترکیبی خطی از ستون‌های $A$ است.}

\paragraph{۴. مجموعه‌ای از ضرب‌های بردار در ماتریس.} سطر $i$-ام $C$ برابر $\va_i\T B$ است:
\[
C=AB=\begin{bmatrix}- & \va_1\T & -\\ - & \va_2\T & -\\ & \vdots & \\ - & \va_m\T & -\end{bmatrix}B=\begin{bmatrix}- & \va_1\T B & -\\ - & \va_2\T B & -\\ & \vdots & \\ - & \va_m\T B & -\end{bmatrix}.
\]
یعنی \hl{هر سطر $AB$ ترکیبی خطی از سطرهای $B$ است.} این نکته کلید فهم عملیات سطری در بخش بعد است.

\subsubsection{ویژگی‌های ضرب ماتریس‌ها}
\begin{theorem}[شرکت‌پذیری و توزیع‌پذیری]\label{thm:assoc}
اگر ابعاد ماتریس‌ها اجازه دهد:
\begin{enumerate}
  \item \textbf{شرکت‌پذیری:} $(AB)C=A(BC)$.
  \item \textbf{توزیع‌پذیری:} $A(B+C)=AB+AC$ و $(A+B)C=AC+BC$.
  \item \textbf{همانی:} اگر $A$ ماتریس $m\times n$ باشد، $I_mA=A=AI_n$.
\end{enumerate}
\end{theorem}
\begin{proof}
\textbf{۱.} فرض کنید $A$ از اندازه‌ی $m\times n$، $B$ از اندازه‌ی $n\times p$ و $C$ از اندازه‌ی $p\times q$ باشد. درایه‌ی $(i,j)$ هر دو طرف را حساب می‌کنیم:
\[
\big((AB)C\big)_{ij}=\sum_{s=1}^{p}(AB)_{is}C_{sj}=\sum_{s=1}^{p}\Big(\sum_{r=1}^{n}A_{ir}B_{rs}\Big)C_{sj}=\sum_{s=1}^{p}\sum_{r=1}^{n}A_{ir}B_{rs}C_{sj},
\]
\[
\big(A(BC)\big)_{ij}=\sum_{r=1}^{n}A_{ir}(BC)_{rj}=\sum_{r=1}^{n}A_{ir}\Big(\sum_{s=1}^{p}B_{rs}C_{sj}\Big)=\sum_{r=1}^{n}\sum_{s=1}^{p}A_{ir}B_{rs}C_{sj}.
\]
هر دو، مجموع همان $np$ جمله‌ی $A_{ir}B_{rs}C_{sj}$ هستند و فقط ترتیب جمع‌زدن فرق دارد؛ پس برابرند.

\textbf{۲.} $\big(A(B+C)\big)_{ij}=\sum_r A_{ir}(B_{rj}+C_{rj})=\sum_r A_{ir}B_{rj}+\sum_r A_{ir}C_{rj}=(AB)_{ij}+(AC)_{ij}$. تساوی دوم هم مشابه است.

\textbf{۳.} $(I_mA)_{ij}=\sum_r\delta_{ir}A_{rj}=A_{ij}$، چون فقط جمله‌ی $r=i$ باقی می‌ماند. مشابهاً $(AI_n)_{ij}=A_{ij}$.
\end{proof}

\begin{warning}[ضرب ماتریس‌ها جابه‌جایی نیست]
در حالت کلی $AB\neq BA$. گاهی حتی ابعاد اجازه‌ی ضرب در یک طرف را نمی‌دهند (مثلاً $A_{2\times3}$ و $B_{3\times 4}$: $AB$ تعریف شده ولی $BA$ نه). حتی برای ماتریس‌های مربعی:
\[
A=\begin{bmatrix}0&1\\0&0\end{bmatrix},\ B=\begin{bmatrix}0&0\\1&0\end{bmatrix}:\qquad
AB=\begin{bmatrix}1&0\\0&0\end{bmatrix}\neq\begin{bmatrix}0&0\\0&1\end{bmatrix}=BA .
\]
برای همین در جبر ماتریسی باید مراقب ترتیب باشیم؛ مثلاً $(A+B)^2=A^2+AB+BA+B^2$ و نه $A^2+2AB+B^2$.
\end{warning}

\begin{theorem}[ترانهاده‌ی حاصل‌ضرب]\label{thm:transpose}
برای هر $A\in\R^{m\times n}$ و $B\in\R^{n\times p}$:
\[
(A\T)\T=A,\qquad (A+B)\T=A\T+B\T\ (\text{\fa{با ابعاد برابر}}),\qquad (AB)\T=B\T A\T .
\]
\end{theorem}
\begin{proof}
دو مورد اول مستقیم از تعریف‌اند. برای سومی، درایه‌ی $(i,j)$ دو طرف را مقایسه می‌کنیم:
\[
\big((AB)\T\big)_{ij}=(AB)_{ji}=\sum_r A_{jr}B_{ri}=\sum_r (B\T)_{ir}(A\T)_{rj}=(B\T A\T)_{ij}.
\]
\end{proof}

\subsection{ماتریس وارون‌پذیر}
\begin{definition}[وارون ماتریس]
ماتریس مربعی $A\in\F^{n\times n}$ را \term{وارون‌پذیر}{Invertible} می‌نامیم اگر ماتریسی مثل $B$ وجود داشته باشد که
\[
AB=BA=I_n .
\]
در این صورت $B$ را \term{وارون}{Inverse} $A$ می‌نامیم و با $A^{-1}$ نشان می‌دهیم.
\end{definition}

\begin{theorem}[ویژگی‌های وارون]\label{thm:inverse-props}
\begin{enumerate}
  \item وارون یک ماتریس (اگر وجود داشته باشد) یکتاست.
  \item اگر $A$ وارون‌پذیر باشد، $A^{-1}$ هم وارون‌پذیر است و $(A^{-1})^{-1}=A$.
  \item اگر $A$ و $B$ (با ابعاد یکسان) وارون‌پذیر باشند، $AB$ هم وارون‌پذیر است و $(AB)^{-1}=B^{-1}A^{-1}$.
  \item اگر $A$ وارون‌پذیر باشد، $A\T$ هم وارون‌پذیر است و $(A\T)^{-1}=(A^{-1})\T$.
\end{enumerate}
\end{theorem}
\begin{proof}
\textbf{۱.} اگر $B$ و $C$ هر دو وارون $A$ باشند، با شرکت‌پذیری:
$B=BI=B(AC)=(BA)C=IC=C$.

\textbf{۲.} تساوی‌های $AA^{-1}=A^{-1}A=I$ دقیقاً می‌گویند $A$ وارون $A^{-1}$ است.

\textbf{۳.} با شرکت‌پذیری:
$(AB)(B^{-1}A^{-1})=A(BB^{-1})A^{-1}=AIA^{-1}=AA^{-1}=I$
و به همین ترتیب $(B^{-1}A^{-1})(AB)=I$. (به عوض شدن ترتیب دقت کنید؛ مثل پوشیدن و درآوردن جوراب و کفش!)

\textbf{۴.} با قضیه‌ی \ref{thm:transpose}: $A\T(A^{-1})\T=(A^{-1}A)\T=I\T=I$ و $(A^{-1})\T A\T=(AA^{-1})\T=I$.
\end{proof}

\begin{note}[وارون ماتریس $2\times2$]
اگر $ad-bc\neq0$ باشد،
\[
\begin{bmatrix}a&b\\c&d\end{bmatrix}^{-1}=\frac{1}{ad-bc}\begin{bmatrix}d&-b\\-c&a\end{bmatrix}.
\]
درستی آن را با ضرب کردن دو ماتریس در هم و رسیدن به $I$ می‌توانید بررسی کنید.
\end{note}

\section{عملیات سطری مقدماتی}

\subsection{سه عمل اصلی}
در «ایده‌ی حذف» دیدیم که با کم کردن مضربی از یک معادله از دیگری، دستگاه ساده‌تری با \emph{همان جواب‌ها} به دست می‌آید. روی ماتریس افزوده این کار یعنی دست بردن در سطرها. سه نوع عمل مجاز داریم:

\begin{definition}[عملیات سطری مقدماتی]
\term{عملیات سطری مقدماتی}{Elementary Row Operations} روی یک ماتریس این سه نوع‌اند:
\begin{enumerate}
  \item \textbf{مقیاس‌دهی} \enterm{Scaling}: ضرب همه‌ی درایه‌های یک سطر در یک اسکالر \emph{ناصفر} $c$. \quad نماد: $R_i\leftarrow cR_i$.
  \item \textbf{جایگزینی} \enterm{Replacement}: جایگزین کردن یک سطر با مجموع خودش و مضربی از سطر دیگر. \quad نماد: $R_i\leftarrow R_i+cR_k$ \ ($k\neq i$).
  \item \textbf{جابه‌جایی} \enterm{Interchange}: عوض کردن جای دو سطر. \quad نماد: $R_i\leftrightarrow R_k$.
\end{enumerate}
\end{definition}

به زبان دقیق‌تر، هر عمل سطری مقدماتی یک \textbf{تابع} $e$ است که به هر ماتریس $m\times n$ مثل $A$ یک ماتریس $m\times n$ مثل $e(A)$ نسبت می‌دهد ($c\neq0$):
\begin{enumerate}
  \item مقیاس‌دهی سطر $r$: \quad $e(A)_{ij}=A_{ij}$ برای $i\neq r$ و $e(A)_{rj}=c\,A_{rj}$.
  \item جایگزینی سطر $r$ با $r+c\times$ سطر $s$: \quad $e(A)_{ij}=A_{ij}$ برای $i\neq r$ و $e(A)_{rj}=A_{rj}+c\,A_{sj}$.
  \item جابه‌جایی سطرهای $r$ و $s$: \quad $e(A)_{ij}=A_{ij}$ برای $i\neq r,s$ و $e(A)_{rj}=A_{sj}$ و $e(A)_{sj}=A_{rj}$.
\end{enumerate}

\begin{note}
در تعریف $e(A)$ مهم نیست $A$ چند ستون دارد؛ اما \textbf{تعداد سطرهای} $A$ اهمیت دارد. مثلاً «جابه‌جایی سطرهای ۴ و ۵» روی یک ماتریس سه‌سطری بی‌معناست.
\end{note}

\begin{warning}
در مقیاس‌دهی، شرط $c\neq0$ ضروری است. اگر سطری را در صفر ضرب کنیم، اطلاعات یک معادله را برای همیشه از دست می‌دهیم و دیگر نمی‌توانیم به عقب برگردیم.
\end{warning}

\subsection{وارون عملیات سطری}
\begin{theorem}[وارون عمل سطری مقدماتی]\label{thm:inverse-rowop}
برای هر عمل سطری مقدماتی $e$، یک عمل سطری مقدماتی $e_1$ \textbf{از همان نوع} وجود دارد که برای هر ماتریس $A$:
\[
e_1\big(e(A)\big)=e\big(e_1(A)\big)=A .
\]
یعنی تابع وارون $e$ وجود دارد و خودش یک عمل سطری مقدماتی از همان نوع است.
\end{theorem}
\begin{proof}
هر سه نوع را جداگانه بررسی می‌کنیم:
\begin{enumerate}
  \item اگر $e$ سطر $r$ را در $c\neq0$ ضرب کند، $e_1$ را ضرب سطر $r$ در $c^{-1}$ بگیرید (این‌جا از ناصفر بودن $c$ استفاده کردیم). سطر $r$ اول در $c$ و بعد در $c^{-1}$ (یا برعکس) ضرب می‌شود و به خودش برمی‌گردد؛ بقیه‌ی سطرها اصلاً تغییر نمی‌کنند.
  \item اگر $e$ سطر $r$ را با $R_r+cR_s$ جایگزین کند، $e_1$ را جایگزینی سطر $r$ با $R_r-cR_s$ بگیرید. چون سطر $s$ در این عمل‌ها تغییر نمی‌کند:
  \[
  (R_r+cR_s)-cR_s=R_r .
  \]
  \item اگر $e$ سطرهای $r$ و $s$ را جابه‌جا کند، خود $e$ را به‌عنوان $e_1$ بگیرید: دو بار جابه‌جا کردن، همه‌چیز را به حالت اول برمی‌گرداند.
\end{enumerate}
\end{proof}

\subsection{هم‌ارزی سطری}
\begin{definition}[هم‌ارز سطری]
اگر $A$ و $B$ ماتریس‌های $m\times n$ روی میدان $\F$ باشند، می‌گوییم $B$ با $A$ \term{هم‌ارز سطری}{Row-Equivalent} است اگر بتوان $B$ را با یک دنباله‌ی \emph{متناهی} از عملیات سطری مقدماتی از $A$ به دست آورد. می‌نویسیم $A\sim B$.
\end{definition}

\begin{theorem}[هم‌ارزی سطری یک رابطه‌ی هم‌ارزی است]\label{thm:row-equiv-relation}
برای ماتریس‌های $m\times n$:
\begin{enumerate}
  \item \textbf{بازتابی:} هر ماتریس با خودش هم‌ارز سطری است.
  \item \textbf{تقارنی:} اگر $B$ با $A$ هم‌ارز سطری باشد، $A$ هم با $B$ هم‌ارز سطری است.
  \item \textbf{تعدی:} اگر $B$ با $A$ و $C$ با $B$ هم‌ارز سطری باشند، $C$ با $A$ هم‌ارز سطری است.
\end{enumerate}
\end{theorem}
\begin{proof}
\textbf{۱.} مثلاً یک سطر را در ۱ ضرب کنید (یا دنباله‌ی تهی از عمل‌ها را در نظر بگیرید).

\textbf{۲.} فرض کنید $B=e_k(\cdots e_2(e_1(A))\cdots)$. طبق قضیه‌ی \ref{thm:inverse-rowop} هر $e_i$ وارونی مثل $e_i'$ دارد که خودش عمل سطری مقدماتی است. حالا عمل‌های وارون را به ترتیب \emph{برعکس} اعمال می‌کنیم:
\[
e_1'\big(e_2'(\cdots e_k'(B)\cdots)\big)=A .
\]
\textbf{۳.} دنباله‌ی عمل‌هایی که $A$ را به $B$ می‌برد را با دنباله‌ای که $B$ را به $C$ می‌برد پشت هم می‌گذاریم؛ دنباله‌ای متناهی به دست می‌آید که $A$ را به $C$ می‌برد.
\end{proof}

به خاطر ویژگی تقارنی، به جای «$B$ با $A$ هم‌ارز سطری است» می‌توانیم بگوییم «$A$ و $B$ هم‌ارز سطری‌اند».

\begin{theorem}[عملیات سطری جواب‌ها را تغییر نمی‌دهند]\label{thm:row-equiv-solutions}
اگر ماتریس‌های افزوده‌ی دو دستگاه خطی هم‌ارز سطری باشند، دو دستگاه دقیقاً مجموعه‌ی جواب یکسانی دارند.
\end{theorem}
\begin{proof}
کافی است برای \emph{یک} عمل سطری مقدماتی ثابت کنیم؛ چون دنباله‌ای از عمل‌ها، قدم‌به‌قدم مجموعه‌ی جواب را حفظ می‌کند. فرض کنید $\vs$ جواب دستگاه اول باشد.
\begin{itemize}
  \item \textbf{مقیاس‌دهی:} اگر معادله‌ی $r$ برای $\vs$ برقرار باشد، $c$ برابر آن هم برقرار است.
  \item \textbf{جایگزینی:} اگر معادله‌های $r$ و $s$ برای $\vs$ برقرار باشند، مجموع معادله‌ی $r$ و $c$ برابر معادله‌ی $s$ هم برقرار است.
  \item \textbf{جابه‌جایی:} فقط ترتیب معادله‌ها عوض شده است.
\end{itemize}
پس هر جواب دستگاه اول، جواب دستگاه جدید است. برعکس، دستگاه اول را می‌توان با عمل وارون (قضیه‌ی \ref{thm:inverse-rowop}) از دستگاه جدید ساخت؛ پس با همین استدلال، هر جواب دستگاه جدید هم جواب دستگاه اول است.
\end{proof}

\section{ماتریس‌های مقدماتی}

\subsection{تعریف}
\begin{definition}[ماتریس مقدماتی]
یک ماتریس $m\times m$ را \term{مقدماتی}{Elementary Matrix} می‌نامیم اگر بتوان آن را با \textbf{دقیقاً یک} عمل سطری مقدماتی از ماتریس همانی $I_m$ به دست آورد.
\end{definition}

\begin{example}[همه‌ی ماتریس‌های مقدماتی $2\times 2$]
همه‌ی ماتریس‌های مقدماتی $2\times2$ را پیدا کنید.
\end{example}
\begin{solution}
سه نوع عمل سطری را روی $I_2=\begin{bmatrix}1&0\\0&1\end{bmatrix}$ اعمال می‌کنیم:
\begin{center}
\begin{tabular}{lcc}
\toprule
\textbf{نوع عمل} & \textbf{عمل سطری} & \textbf{ماتریس مقدماتی}\\
\midrule
مقیاس‌دهی ($c\neq0$) & $R_1\leftarrow cR_1$ & $\begin{bmatrix}c&0\\0&1\end{bmatrix}$\\[5mm]
مقیاس‌دهی ($c\neq0$) & $R_2\leftarrow cR_2$ & $\begin{bmatrix}1&0\\0&c\end{bmatrix}$\\[5mm]
جایگزینی ($c$ دلخواه) & $R_1\leftarrow R_1+cR_2$ & $\begin{bmatrix}1&c\\0&1\end{bmatrix}$\\[5mm]
جایگزینی ($c$ دلخواه) & $R_2\leftarrow R_2+cR_1$ & $\begin{bmatrix}1&0\\c&1\end{bmatrix}$\\[5mm]
جابه‌جایی & $R_1\leftrightarrow R_2$ & $\begin{bmatrix}0&1\\1&0\end{bmatrix}$\\
\bottomrule
\end{tabular}
\end{center}
\end{solution}

\subsection{عملیات سطری = ضرب از چپ در ماتریس مقدماتی}
نکته‌ی جادویی این است که \emph{هر عمل سطری را می‌توان با یک ضرب ماتریسی انجام داد.}

\begin{theorem}[ماتریس مقدماتی و عمل سطری]\label{thm:elem-matrix}
فرض کنید $e$ یک عمل سطری مقدماتی و $E=e(I_m)$ ماتریس مقدماتی متناظر با آن باشد. آن‌گاه برای \textbf{هر} ماتریس $m\times n$ مثل $A$:
\[
e(A)=EA .
\]
\end{theorem}
\begin{proof}
درایه‌ی $(i,j)$ حاصل‌ضرب برابر است با $(EA)_{ij}=\sum_{k=1}^m E_{ik}A_{kj}$. سطرهای $E$ را برای هر سه نوع عمل می‌شناسیم، چون $E$ همان $I$ است که یک عمل رویش انجام شده (یادآوری: $I_{ik}=\delta_{ik}$).

\textbf{نوع ۱ (ضرب سطر $r$ در $c$):} برای $i\neq r$ داریم $E_{ik}=\delta_{ik}$ و $E_{rk}=c\,\delta_{rk}$. پس
\[
(EA)_{ij}=\sum_k\delta_{ik}A_{kj}=A_{ij}\ \ (i\neq r),\qquad (EA)_{rj}=\sum_k c\,\delta_{rk}A_{kj}=c\,A_{rj}.
\]
\textbf{نوع ۲ (سطر $r$ با $R_r+cR_s$):} برای $i\neq r$ داریم $E_{ik}=\delta_{ik}$ و $E_{rk}=\delta_{rk}+c\,\delta_{sk}$. پس
\[
(EA)_{ij}=A_{ij}\ \ (i\neq r),\qquad (EA)_{rj}=\sum_k(\delta_{rk}+c\,\delta_{sk})A_{kj}=A_{rj}+c\,A_{sj}.
\]
\textbf{نوع ۳ (جابه‌جایی $r$ و $s$):} برای $i\neq r,s$ داریم $E_{ik}=\delta_{ik}$، و $E_{rk}=\delta_{sk}$، $E_{sk}=\delta_{rk}$. پس
\[
(EA)_{ij}=A_{ij}\ \ (i\neq r,s),\qquad (EA)_{rj}=A_{sj},\qquad (EA)_{sj}=A_{rj}.
\]
در هر سه حالت $EA$ دقیقاً همان $e(A)$ است.
\end{proof}

\begin{intuition}
از نگاه چهارم ضرب ماتریس‌ها می‌دانیم \textbf{هر سطر $EA$ ترکیبی خطی از سطرهای $A$ است} و ضرایب این ترکیب را سطر متناظر $E$ تعیین می‌کند. پس «ضرب از چپ در یک ماتریس مربعی» یعنی «ساختن سطرهای جدید از روی سطرهای قدیمی». به زبان ساده:
\begin{center}
\hl{ضرب از چپ در یک ماتریس مربعی، عملیات سطری انجام می‌دهد.}
\end{center}
\end{intuition}

\begin{example}[تبدیل یک ماتریس به $I$ با ماتریس‌های مقدماتی]
ماتریس $A=\begin{bmatrix}1&0&2\\-2&0&-3\\0&2&0\end{bmatrix}$ را با عملیات سطری به $I$ تبدیل کنید و ماتریس مقدماتی هر مرحله را بنویسید.
\end{example}
\needspace{18\baselineskip}
\begin{solution}
\begin{center}
\begin{tabular}{ccc}
\toprule
\textbf{ماتریس پس از عمل} & \textbf{عمل سطری} & \textbf{ماتریس مقدماتی} \\
\midrule
$\begin{bmatrix}1&0&2\\0&0&1\\0&2&0\end{bmatrix}$ & $R_2\leftarrow R_2+2R_1$ & $E_1=\begin{bmatrix}1&0&0\\2&1&0\\0&0&1\end{bmatrix}$ \\[6mm]
$\begin{bmatrix}1&0&2\\0&2&0\\0&0&1\end{bmatrix}$ & $R_2\leftrightarrow R_3$ & $E_2=\begin{bmatrix}1&0&0\\0&0&1\\0&1&0\end{bmatrix}$ \\[6mm]
$\begin{bmatrix}1&0&2\\0&1&0\\0&0&1\end{bmatrix}$ & $R_2\leftarrow \frac12 R_2$ & $E_3=\begin{bmatrix}1&0&0\\0&\frac12&0\\0&0&1\end{bmatrix}$ \\[6mm]
$\begin{bmatrix}1&0&0\\0&1&0\\0&0&1\end{bmatrix}$ & $R_1\leftarrow R_1-2R_3$ & $E_4=\begin{bmatrix}1&0&-2\\0&1&0\\0&0&1\end{bmatrix}$ \\
\bottomrule
\end{tabular}
\end{center}
طبق قضیه‌ی \ref{thm:elem-matrix} و شرکت‌پذیری ضرب:
\[
E_4\big(E_3\big(E_2(E_1A)\big)\big)=(E_4E_3E_2E_1)A=I .
\]
پس ماتریس $P=E_4E_3E_2E_1$ همه‌ی کار را یک‌جا انجام می‌دهد. (در فصل بعد می‌بینیم که این یعنی $P=A^{-1}$.)
\end{solution}

\subsection{هم‌ارزی سطری و ماتریس‌های مقدماتی}
\begin{theorem}\label{thm:row-equiv-PA}
فرض کنید $A$ و $B$ ماتریس‌های $m\times n$ باشند. $B$ با $A$ هم‌ارز سطری است اگر و تنها اگر
\[
B=PA,
\]
که در آن $P$ حاصل‌ضرب چند ماتریس مقدماتی $m\times m$ است.
\end{theorem}
\begin{proof}
\textbf{(رفت)} فرض کنید $B$ با عمل‌های $e_1,\dots,e_k$ (به همین ترتیب) از $A$ به دست آمده باشد و $E_i=e_i(I)$. با $k$ بار استفاده از قضیه‌ی \ref{thm:elem-matrix}:
\[
B=e_k(\cdots e_2(e_1(A))\cdots)=E_k(\cdots(E_2(E_1A))\cdots)=(E_k\cdots E_2E_1)A,
\]
که در قدم آخر از شرکت‌پذیری استفاده کردیم. پس $P=E_k\cdots E_1$.

\textbf{(برگشت)} فرض کنید $B=E_k\cdots E_2E_1A$ و هر $E_i$ مقدماتی باشد؛ یعنی $E_i=e_i(I)$ برای یک عمل سطری $e_i$. باز طبق قضیه‌ی \ref{thm:elem-matrix}، $E_1A=e_1(A)$، سپس $E_2(E_1A)=e_2(e_1(A))$ و الی آخر. پس $B$ با دنباله‌ی $e_1,\dots,e_k$ از $A$ به دست می‌آید.
\end{proof}

\begin{corollary}[ماتریس‌های مقدماتی وارون‌پذیرند]\label{cor:elem-invertible}
هر ماتریس مقدماتی $E$ وارون‌پذیر است و وارونش یک ماتریس مقدماتی از همان نوع است. در نتیجه هر حاصل‌ضرب از ماتریس‌های مقدماتی (مثل $P$ در قضیه‌ی قبل) وارون‌پذیر است.
\end{corollary}
\begin{proof}
فرض کنید $E=e(I)$ و $e_1$ عمل وارون $e$ باشد (قضیه‌ی \ref{thm:inverse-rowop}) و $E_1=e_1(I)$. طبق قضیه‌ی \ref{thm:elem-matrix}:
\[
E_1E=e_1(E)=e_1(e(I))=I,\qquad EE_1=e(E_1)=e(e_1(I))=I .
\]
پس $E^{-1}=E_1$. بخش دوم از قضیه‌ی \ref{thm:inverse-props} (وارون حاصل‌ضرب) نتیجه می‌شود.
\end{proof}

مثلاً وارون $\begin{bmatrix}1&0&0\\2&1&0\\0&0&1\end{bmatrix}$ ($R_2\leftarrow R_2+2R_1$) برابر $\begin{bmatrix}1&0&0\\-2&1&0\\0&0&1\end{bmatrix}$ ($R_2\leftarrow R_2-2R_1$) است.

\subsection{یک مثال روی اعداد مختلط}
\begin{example}[دستگاه همگن روی $\C$]
فرض کنید میدان، اعداد مختلط باشد و ماتریس ضرایب یک دستگاه همگن این باشد:
\[
A=\begin{bmatrix}-1&i\\-i&3\\1&2\end{bmatrix}.
\]
همه‌ی جواب‌های $A\vx=\zero$ را پیدا کنید.
\end{example}
\begin{solution}
برای دستگاه همگن لازم نیست ستون صفرِ سمت راست را بنویسیم، چون عملیات سطری آن را صفر نگه می‌دارند. با عملیات سطری (و استفاده از $i^2=-1$):
\[
\begin{bmatrix}-1&i\\-i&3\\1&2\end{bmatrix}
\xrightarrow[R_3\leftarrow R_3+R_1]{R_2\leftarrow R_2-iR_1}
\begin{bmatrix}-1&i\\0&3-i^2\\0&2+i\end{bmatrix}
=\begin{bmatrix}-1&i\\0&4\\0&2+i\end{bmatrix}
\xrightarrow{R_2\leftarrow\frac14R_2}
\begin{bmatrix}-1&i\\0&1\\0&2+i\end{bmatrix}
\]
\[
\xrightarrow[R_3\leftarrow R_3-(2+i)R_2]{R_1\leftarrow R_1-iR_2}
\begin{bmatrix}-1&0\\0&1\\0&0\end{bmatrix}
\xrightarrow{R_1\leftarrow -R_1}
\begin{bmatrix}1&0\\0&1\\0&0\end{bmatrix}.
\]
(در قدم اول: درایه‌ی دوم سطر دوم برابر $3-i\cdot i=3-i^2=4$ است.) دستگاه نهایی می‌گوید $x_1=0$ و $x_2=0$. طبق قضیه‌ی \ref{thm:row-equiv-solutions} دستگاه اولیه هم همین جواب را دارد: \textbf{فقط جواب بدیهی} $\vx=\zero$.
\end{solution}

\section{جمع‌بندی و تمرین}

\begin{summary}
\begin{itemize}
  \item هر دستگاه خطی را می‌توان به شکل $A\vx=\vb$ نوشت؛ ماتریس افزوده $\aug{A}{\vb}$ همه‌ی اطلاعات آن را دارد.
  \item دستگاه خطی یا جواب ندارد، یا یک جواب دارد، یا بی‌شمار جواب.
  \item ضرب داخلی $\va\T\vb=\sum a_ib_i$ یک عدد است؛ $\norm{\va}^2=\va\T\va$ و $\va\cdot\vb=\norm{\va}\norm{\vb}\cos\theta$.
  \item $A\vx$ ترکیب خطی ستون‌های $A$ است و $\vx\T A$ ترکیب خطی سطرهای $A$.
  \item ضرب ماتریس‌ها شرکت‌پذیر و توزیع‌پذیر است اما جابه‌جایی نیست؛ $(AB)\T=B\T A\T$ و $(AB)^{-1}=B^{-1}A^{-1}$.
  \item سه عمل سطری مقدماتی (مقیاس‌دهی، جایگزینی، جابه‌جایی) وارون‌پذیرند و جواب دستگاه را تغییر نمی‌دهند.
  \item هر عمل سطری همان ضرب از چپ در یک ماتریس مقدماتی است: $e(A)=EA$. پس $A\sim B\iff B=PA$ که $P$ حاصل‌ضرب ماتریس‌های مقدماتی (و وارون‌پذیر) است.
\end{itemize}
\end{summary}

\begin{exercise}
برای $A=\begin{bmatrix}1&2\\0&1\end{bmatrix}$ و $B=\begin{bmatrix}1&0\\3&1\end{bmatrix}$، حاصل $AB$ و $BA$ را حساب کنید.
\end{exercise}

\begin{exercise}
نشان دهید اگر $AB=0$ باشد، لزوماً $A=0$ یا $B=0$ نیست.
\end{exercise}

\begin{exercise}
زاویه‌ی بین $\vu=(1,1,0)$ و $\vv=(1,0,1)$ را پیدا کنید.
\end{exercise}

\begin{exercise}
ماتریس مقدماتی $3\times3$ را بنویسید که عمل $R_3\leftarrow R_3-4R_1$ را انجام می‌دهد. وارون آن چیست؟
\end{exercise}

\begin{exercise}
ثابت کنید اگر برای \emph{هر} $\vx\in\R^n$ داشته باشیم $\va\cdot\vx=0$، آن‌گاه $\va=\zero$.
\end{exercise}

\begin{exercise}
نشان دهید اگر $A$ وارون‌پذیر باشد، دستگاه $A\vx=\vb$ برای هر $\vb$ دقیقاً یک جواب دارد.
\end{exercise}

\subsection*{پاسخ‌های کوتاه}
\begin{enumerate}
  \item $AB=\begin{bmatrix}7&2\\3&1\end{bmatrix}$ و $BA=\begin{bmatrix}1&2\\3&7\end{bmatrix}$؛ پس $AB\neq BA$.
  \item $A=B=\begin{bmatrix}0&1\\0&0\end{bmatrix}$: هر دو ناصفرند اما $AB=0$.
  \item $\cos\theta=\frac{\vu\cdot\vv}{\norm{\vu}\norm{\vv}}=\frac{1}{\sqrt2\sqrt2}=\frac12$، پس $\theta=60^\circ$.
  \item $E=\begin{bmatrix}1&0&0\\0&1&0\\-4&0&1\end{bmatrix}$ و $E^{-1}=\begin{bmatrix}1&0&0\\0&1&0\\4&0&1\end{bmatrix}$.
  \item $\vx=\va$ را بگذارید: $\va\cdot\va=0$ و طبق ویژگی مثبت بودن $\va=\zero$.
  \item \textbf{وجود:} $\vx=A^{-1}\vb$ جواب است چون $A(A^{-1}\vb)=\vb$. \textbf{یکتایی:} اگر $A\vx=\vb$، دو طرف را از چپ در $A^{-1}$ ضرب کنید: $\vx=A^{-1}\vb$.
\end{enumerate}
